Now, a human being really finds in himself a capacity by which he distinguishes
himself from all other things, even from himself insofar as he is affected by objects,
and that is reason. This, as pure self activity, is raised even above understanding by this:
that though the latter is also self activity and does not, like sense, contain merely representations
that arise when we are affected by things (and are thus passive),
yet it can produce from its activity no other concepts
than those which serve merely to bring sensible representations under rules
and thereby to unite them in one consciousness,
without which use of sensibility it would think nothing at all;
but reason, on the contrary shows, in what we call 'ideas' a spontaneity
so pure that it thereby goes far beyond anything that sensibility
can ever afford it, and proves its highest occupation in distinguishing
the world of sense and the world of understanding  from each other
and thereby marking out limits for the understanding itself.

